\section{Two Hecke correspondences and a quadratic test}
\label{sec:twohecke}

We now convert the Hecke vanishing of Lemma~\ref{spec:vanishing}
into a local test for a non-nilpotent cotangent direction.
The first Hecke correspondence supplies a proper pushforward.
The second makes its pullback vanish, while the two transverse
sections of Proposition~\ref{hecke:transverse} isolate one stalk
and identify its neighborhood with a split quadratic equation.
The next section will use a projective quadric to remove the
extra variables in this equation.

\begin{proposition}[Vanishing at the quadratic vertex]
\label{two:vertex}
Let $F\in\Shv(\mathcal B_K)$ be the object of
Proposition~\ref{spec:missing}, with the vanishings of
Lemma~\ref{spec:vanishing}. Let $b:U\to\mathcal B$ be a smooth
chart, where $U$ is a scheme of finite type over $S$.
Let $f\in\mathcal O(U)$ and $u\in U_s(k)$ satisfy $f(u)=0$.
Suppose
\[
 df_u=b^*A
 \quad\text{for a non-nilpotent }A\in T^*_{b(u)}\mathcal B_s.
\]
Then there exist $x\in\mathscr X_s(k)$ and an integer $m\geq0$
with the following property. Put $T=U\times_S\mathscr X$,
let $F_T$ be the pullback of $F_U=b_K^*F$, and set
\[
 Q^a=\left\{\sum_{i=1}^m u_i v_i+f=0\right\}
       \subset T\times_S\mathbb A_S^{2m},
 \qquad \tau:Q^a\longrightarrow T.
\]
Here $u_i,v_i$ are affine coordinates, distinct from the point $u$,
and $f$ is pulled back from $U$. At the point
$\nu=(u,x,0,0)$ one has
\begin{equation}
 \label{eq:two-vertex}
 \bigl(\Psi_{Q^a}(\tau_K^*F_T)\bigr)_\nu=0.
\end{equation}
When $m=0$, this is the vanishing on
$Q^a=V(f)\times_S\mathscr X$.
\end{proposition}

\begin{proof}
Choose the modification $h$, its leg $x$, its orbit dimension $m$,
and $\theta=(A',\xi)$ from Proposition~\ref{hecke:transverse}.
Thus
\begin{equation}
 \label{eq:two-cotangent}
 p^*A=q^*\theta\quad\text{at }h,
 \qquad \xi\ne0.
\end{equation}
We use the exact-position correspondence $H=H^\lambda$, its proper
bound $\overline H$, and its generic Satake integrand
$\mathcal I_K$ from Lemma~\ref{hecke:models}.
On $H_K$ this integrand is $p_K^*F$ tensored with a fixed invertible
shift and Tate line. Such a factor does not affect vanishing;
we suppress it after restricting to the open orbit.

\paragraph{A proper fiber with zero cohomology.}
Write $D=V(f)$. Take two copies $H_1,H_2$ of $H$ and form
\[
 W_0=H_2\times_{p_2,\mathcal B,b}U,
 \qquad W=H_2\times_{\mathcal B}D\subset W_0.
\]
Both map to $\mathcal Z=\mathcal B\times_S\mathscr X$ through
$q_2=(o_2,\ell_{H_2})$. The map $W_0\to\mathcal Z$ is smooth:
its projection to $H_2$ is the base change of $b$, and $q_2$ is
smooth. At $w=(h,u)$ the differential of the equation cutting
out $W$ is the pullback of $\theta$ by
\eqref{eq:two-cotangent}. Its relative differential over the
output bundle stack is nonzero, because the leg component is
$\xi\ne0$. The smooth hypersurface criterion therefore makes
$W\to\mathcal B$ by the output bundle smooth near $w$.
The arbitrary-leg assertion of Lemma~\ref{spec:vanishing} gives
\begin{equation}
 \label{eq:two-hecke-zero}
 \bigl(\Psi_W(q_{2,K}^*\mathrm H_V(F))\bigr)_w=0.
\end{equation}
This uses smoothness only of the projection to the bundle stack.

The two fiber products used below are
\[
 P=H_1\times_{\mathcal Z}W_0,
 \qquad
 P_D=H_1\times_{\mathcal Z}W
       \subset\overline P_D=\overline H_1\times_{\mathcal Z}W.
\]
The following squares are cartesian; the second vertical map
$\pi$ is proper. In the first square, $\Delta$ repeats the
modification in the two factors.
\[
\begin{tikzcd}[column sep=large,row sep=large]
 P \arrow[r] \arrow[d] & H_1 \arrow[d,"q_1"] \\
 W_0 \arrow[r,"q_2"'] \arrow[u,bend left=38,"\Delta"]
       & \mathcal Z
\end{tikzcd}
\qquad
\begin{tikzcd}[column sep=large,row sep=large]
 \overline P_D \arrow[r] \arrow[d,"\pi"']
       & \overline H_1 \arrow[d,"\overline q_1"] \\
 W \arrow[r,"q_2"'] & \mathcal Z.
\end{tikzcd}
\]
Write $c:P\to U$ for the map through $W_0$, and
$a:P\to\mathcal B$ for the input map $p_1$.
Thus $a$ records the first input, from which the sheaf is pulled back,
while $c$ records the second input in $U$, on which we impose $f=0$.
The distinguished point is $z=(h,h,u)\in P_D(k)$.
All these fiber products are algebraic spaces, because the
fixed-side Hecke maps are representable.

Let $\overline{\mathcal I}_K$ be the pullback of the Satake
integrand to $(\overline P_D)_K$. Proper base change on the
generic fiber identifies
\[
 R\pi_{K,*}\overline{\mathcal I}_K
       \simeq q_{2,K}^*\mathrm H_V(F).
\]
Apply proper compatibility of specialization and then proper
base change at $w$ to \eqref{eq:two-hecke-zero}. With
\[
 Z_w=\pi^{-1}(w),\qquad
 \mathcal K=
   \bigl(\Psi_{\overline P_D}\overline{\mathcal I}_K\bigr)|_{Z_w},
\]
we obtain
\begin{equation}
 \label{eq:two-fiber-zero}
 R\Gamma(Z_w,\mathcal K)=0.
\end{equation}
The space $Z_w$ is the proper bounded Hecke fiber with output
bundle and leg fixed, including their identifications.

\paragraph{Isolating the distinguished stalk.}
The map $a:P\to\mathcal B$ is smooth: $P\to H_1$ is the base
change of $W_0\to\mathcal Z$, and $p_1$ is smooth.
At a point $(h_1,w)$ in the exact-position part of $Z_w$, the
differential of $f(c)$ is the pullback of $q_1^*\theta$.
If its relative differential over $\mathcal B$ is zero, smooth
cotangent injectivity for $P\to H_1$ implies that
$q_1^*\theta$ comes from the input cotangent under $p_1$.
Pass further to the relative cotangent for
$\widetilde p_1=(p_1,\ell_{H_1})$. The leg term disappears,
so the resulting necessary condition is
\[
 S_{A'}(h_1)=0.
\]
By Proposition~\ref{hecke:transverse}, this section has an isolated
zero at $h$ on the fixed-output-and-leg fiber. Consequently
$P_D\to\mathcal B$ by $a$ is smooth at every point of a punctured
neighborhood of $z$ in $Z_w$.
On the exact-position open, the generic integrand is $a_K^*F$.
Lemma~\ref{spec:vanishing} therefore gives an open neighborhood
$O$ of $z$ in $Z_w$ such that
\[
 \mathcal K|_{O\setminus\{z\}}=0.
\]

This local vanishing prevents cancellation with the boundary of
the Hecke bound. Indeed, put $C_w=Z_w\setminus O$ with its reduced
structure. The closed subsets $\{z\}$ and $C_w$ are disjoint.
Open--closed localization identifies $\mathcal K$ with its
restriction to their union, extended by zero. Hence
\[
 \mathcal K\simeq
   i_{z,*}\mathcal K_z\oplus i_{C_w,*}i_{C_w}^*\mathcal K,
 \qquad
 R\Gamma(Z_w,\mathcal K)\simeq
   \mathcal K_z\oplus R\Gamma(C_w,i_{C_w}^*\mathcal K).
\]
The residue field of $z$ is algebraically closed, so the first
cohomology summand is exactly its stalk complex.
Equation~\eqref{eq:two-fiber-zero} forces that summand to vanish.
Removing the invertible Satake normalization gives
\begin{equation}
 \label{eq:two-isolated-zero}
 \bigl(\Psi_{P_D}(a_K^*F)\bigr)_z=0.
\end{equation}
No formula for the Satake kernel on the boundary has been used.

\paragraph{A chart lift along the whole diagonal.}
To interpret \eqref{eq:two-isolated-zero}, we next identify both
the local space and its sheaf. The diagonal
$\Delta\simeq W_0\subset P$ is a section of the separated smooth
map $P\to W_0$ of relative dimension $m$.
It is therefore a closed regular immersion near $z$.
On $\Delta$ there is a specified isomorphism $bc\simeq a$,
coming from the two identical modifications.

We claim that, after an \emph{étale} replacement of $P$ near $z$,
there is a map $r:P\to U$ with an isomorphism
\begin{equation}
 \label{eq:two-chart-lift}
 br\simeq a,
 \qquad r|_\Delta=c|_\Delta,
\end{equation}
whose restriction to $\Delta$ is the specified isomorphism.
For this, form the smooth space
$E_P=P\times_{a,\mathcal B,b}U\to P$.
The given diagonal data define a section over $\Delta$.
First take an étale scheme neighborhood in $P$.
Since $\mathcal B$ has affine diagonal by the representability
statement in Section~\ref{sec:specialization}, $E_P\to P\times_S U$
is affine, so $E_P$ is a scheme as well.
Near the section point, smooth coordinates give an étale map
$E_P\to\mathbb A_P^e$
\cite[Tag~039P]{Stacks}.
After shrinking, the section over the closed diagonal lands in
this coordinate neighborhood. Extend its coordinate functions
from $\Delta$ to $P$, and pull back the étale map along the
extended tuple. This gives an étale neighborhood of $P$ with
a map to $E_P$. Its inverse image of $\Delta$ contains the
open branch specified by the original section.
Remove the closed complement of this branch in that inverse
image. The retained diagonal is the prescribed one, and the
map to $E_P$ proves \eqref{eq:two-chart-lift}, including its
isomorphism of bundles. We retain the notation $P,\Delta,z$.

\paragraph{Exact quadratic coordinates.}
After shrinking $U$, suppose that $U/S$ has constant relative
dimension $N$. The map
\[
 \Delta\longrightarrow T=U\times_S\mathscr X,
 \qquad (c,\text{leg}),
\]
is smooth of relative dimension $m$. Thus $T$, $\Delta$, and
$P$ are smooth over $S$ of relative dimensions
\[
 N+1,\qquad N+1+m,\qquad N+1+2m,
\]
respectively. Choose generators $v_1,\ldots,v_m$ of the
diagonal ideal whose classes form a conormal frame.
Under the identification with
$\Omega_{P/W_0}|_\Delta$, the classes $dv_i$ form a basis.
Because $r=c$ on $\Delta$, there exist regular functions
$u_1,\ldots,u_m$ with the \emph{exact} identity
\begin{equation}
 \label{eq:two-exact-quadratic}
 f(c)-f(r)=\sum_{i=1}^m u_i v_i.
\end{equation}
The restrictions $u_i|_\Delta$ are determined by the conormal
class of the left side.

Consider the entire local diagonal fiber
$\Delta_{(u,x)}$ on the special fiber.
It parametrizes modifications with the input bundle
$b(u)$, its identification, and the leg $x$ fixed.
The relative differential of $f(c)$ for $P\to W_0$ is zero.
On this diagonal fiber, the relative differential of $f(r)$
is induced by $df_u=b^*A$: the restriction $r|_{\Delta_{(u,x)}}$
is constant with value $u$, and the bundle isomorphism in
\eqref{eq:two-chart-lift} is the canonical one along the whole
diagonal.
The conormal coefficient section in
\eqref{eq:two-exact-quadratic} is therefore
\begin{equation}
 \label{eq:two-coefficients}
 \sum_i(u_i|_\Delta)[dv_i]=-S_A
       \quad\text{on }\Delta_{(u,x)}.
\end{equation}
In particular, $u_i(z)=0$, and the derivative of the tuple
$(u_i|_\Delta)$ is an isomorphism on
$T_z\Delta_{(u,x)}$, by Proposition~\ref{hecke:transverse}.
The equality on the whole fiber in
\eqref{eq:two-coefficients} is what supplies this derivative;
equality only at $z$ would not suffice.

Now define
\[
 \Phi=(r,\text{leg},u_1,\ldots,u_m,v_1,\ldots,v_m):
 P\longrightarrow T\times_S\mathbb A_S^{2m}.
\]
Its relative differential at $z$ is an isomorphism.
Indeed, a tangent vector in its kernel is first killed by every
$dv_i$, so lies in $T_z\Delta$. Its $(dr,d\text{leg})$ component
then places it in $T_z\Delta_{(u,x)}$. Finally the derivative
isomorphism in \eqref{eq:two-coefficients} forces the vector
to be zero. The source and target are smooth over $S$ with
the same relative dimension $N+1+2m$, so this injective
differential is an isomorphism. The relative Jacobian
criterion makes $\Phi$ étale at $z$.

By \eqref{eq:two-exact-quadratic}, the cut $P_D=V(f(c))$ is
the exact inverse image of $Q^a$ under $\Phi$. Therefore the
induced map $\Phi_D:P_D\to Q^a$ is étale near $z$, taking $z$ to $\nu$.
Moreover, \eqref{eq:two-chart-lift} identifies
\[
 a_K^*F\simeq r_K^*F_U
       \simeq \Phi_{D,K}^*\tau_K^*F_T
       \quad\text{on }(P_D)_K.
\]
Thus this coordinate map identifies the actual generic sheaf
as well as the equation. Étale base change transforms
\eqref{eq:two-isolated-zero} into \eqref{eq:two-vertex}.

If $m=0$, the diagonal is open near $z$, the sum in
\eqref{eq:two-exact-quadratic} is empty, and the same argument
gives $Q^a=D\times_S\mathscr X$. In that case smooth base
change along $D\times_S\mathscr X\to D$ already yields
$(\Psi_D(F_U|_{D_K}))_u=0$.
\end{proof}

The argument produced exact étale coordinates over the trait.
It used no division by $2$ and no assertion about singular
support in mixed characteristic. Its output for $m>0$ is a
vanishing on the quadratic test space; we next recover the
vanishing on $D$ itself.
